\chapter{Introducción}

\section{Motivación}

En los últimos años, la preocupación por la seguridad en el desarrollo
de software ha crecido notablemente. Como respuesta, los compiladores
modernos han incorporado mecanismos de protección que refuerzan la
integridad y robustez del programa generado.
Estas medidas, sin embargo, introducen un coste en términos de
eficiencia. Se observa un impacto en el tiempo de ejecución, en el uso
de memoria y, en general, en el rendimiento del programa.

Por ello, resulta de interés analizar en qué medida estas protecciones
afectan al comportamiento de los programas, especialmente cuando se
combinan con distintos niveles de optimización. En este trabajo se
estudia el efecto de activar \emph{flags} de seguridad junto con
\emph{flags} de optimización durante el proceso de compilación,
observando cómo esta combinación influye en el rendimiento final.

El análisis se realiza evaluando diferentes compiladores, lo que permite
identificar posibles diferencias en las estrategias utilizadas por estos 
compiladores para aplicar estas medidas. Este enfoque busca comprender 
con mayor precisión el equilibrio entre seguridad y eficiencia, aportando 
datos que ayuden a fundamentar decisiones técnicas en entornos donde ambos
factores son críticos.

\section{Objetivos}

El objetivo principal de este trabajo es realizar un análisis
comparativo del impacto que tienen ciertas medidas de seguridad
aplicadas durante la compilación sobre la eficiencia de ejecución de los
programas generados.
Para ello, se han evaluado compiladores de uso común, como \emph{GCC},
\emph{Clang} y \emph{Rustc} y se han diseñado pruebas en las que se
combinan distintas medidas de seguridad con niveles de optimización.

Durante la experimentación, se recogen métricas como el tiempo de
ejecución, el uso de memoria, el tamaño del programa generado y la
presencia de funciones fortificadas (son aquellas que sustituyen
fragmentos de código potencialmente peligrosos por versiones más seguras
con comprobaciones adicionales). Estas métricas permiten evaluar el
efecto real de la combinación entre seguridad y optimización en cada
entorno de compilación.

Con los resultados obtenidos, se pretende entender cómo responde cada
compilador ante estas configuraciones y en qué medida se ve afectado el
rendimiento. Se busca así ofrecer una visión clara y práctica sobre el
coste asociado a la seguridad en el contexto de la compilación de
programas.

\section{Estructura del documento}

Este documento se encuentra dividido en seis capítulos y dos anexos. En el 
\hyperref[chap:fundamentos]{\textbf{Capítulo 2}} se presentan los conceptos previos necesarios para comprender el 
análisis, con especial atención a las protecciones de seguridad aplicadas a nivel de compilador. En el 
\hyperref[chap:metodologia]{\textbf{Capítulo 3}} se describe el \emph{pipeline} automatizado desarrollado para realizar 
los experimentos: desde la entrada (código fuente y combinaciones de \emph{flags}), pasando por el proceso 
(compilación, ejecución, recolección de métricas), hasta la generación automática de informes y gráficas. También 
se justifican los programas utilizados y el enfoque visual adoptado. En el 
\hyperref[chap:resultados]{\textbf{Capítulo 4}} se detallan el entorno de experimentación, los conjuntos de datos 
utilizados y los resultados obtenidos, incluyendo una discusión crítica sobre las observaciones realizadas. En el
\hyperref[chap:trabajo]{\textbf{Capítulo 5}} se detallan los trabajos similares al que se ha realizado.
Finalmente, en el \hyperref[chap:conclusiones]{\textbf{Capítulo 6}} se exponen las conclusiones obtenidas  tras la realización
de este trabajo, así como posibles líneas de trabajo futuro que pueden continuar y ampliar el estudio aquí presentado.
Adicionalmente, en el \hyperref[chap:anexoA]{\textbf{Anexo A}} se incluye el diagrama de Gantt con la distribución temporal del trabajo, en
el \hyperref[chap:anexoB]{\textbf{Anexo B}} se encuentran la colección completa de gráficas generadas, que no han sido incorporadas en el cuerpo principal del documento.